\section{Иллюстрации результатов}

Ниже приведены основные визуализации, полученные в ходе выполнения лабораторной работы. Интерактивные версии графиков находятся в каталоге \texttt{outputs/}.

\begin{figure}[H]
\centering
\includegraphics[width=\textwidth]{01_target_function.png}
\caption{Функция варианта~3 и обучающие точки с шагом 0,1.}
\label{fig:target}
\end{figure}

\begin{figure}[H]
\centering
\includegraphics[width=\textwidth]{02_training_animation.png}
\caption{Мультипликация обучения traingd ($S = 10$): финальный кадр. Интерактивная версия с кнопками <<Старт>>, <<Пауза>>, <<Сброс>>~--- \texttt{outputs/02\_training\_animation.html}.}
\label{fig:anim}
\end{figure}

\begin{figure}[H]
\centering
\includegraphics[width=\textwidth]{03_training_snapshots.png}
\caption{Изгибание кривой аппроксимации в ходе обучения traingd ($S = 10$).}
\label{fig:snapshots}
\end{figure}

\begin{figure}[H]
\centering
\includegraphics[width=\textwidth]{04_series_approximations.png}
\caption{Аппроксимация при $S = 1, 2, 5, 10, 20, 30$ (traingd, $\eta = 0{,}05$, 600 эпох, типичный запуск).}
\label{fig:series}
\end{figure}

\begin{figure}[H]
\centering
\includegraphics[width=\textwidth]{05_series_performance.png}
\caption{Кривые ошибки traingd при разном числе нейронов.}
\label{fig:seriesperf}
\end{figure}

\begin{figure}[H]
\centering
\includegraphics[width=\textwidth]{06_series_mse_spread.png}
\caption{Разброс итоговой MSE по 10 инициализациям.}
\label{fig:spread}
\end{figure}

\begin{figure}[H]
\centering
\includegraphics[width=\textwidth]{07_algorithms_convergence.png}
\caption{Скорость сходимости traingd, traingdx и trainlm при $S = 10$.}
\label{fig:algos}
\end{figure}

\begin{figure}[H]
\centering
\includegraphics[width=\textwidth]{08_capacity_curve.png}
\caption{Точность аппроксимации в зависимости от числа нейронов скрытого слоя.}
\label{fig:capacity}
\end{figure}

\begin{figure}[H]
\centering
\includegraphics[width=\textwidth]{09_well_trained_models.png}
\caption{Сети, обученные trainlm до сходимости.}
\label{fig:welltrained}
\end{figure}

\begin{figure}[H]
\centering
\includegraphics[width=\textwidth]{10_neuron_contributions.png}
\caption{Выход сети как сумма вкладов нейронов скрытого слоя ($S = 3$ и $S = 5$).}
\label{fig:contrib}
\end{figure}

\begin{figure}[H]
\centering
\includegraphics[width=\textwidth]{11_learning_rate.png}
\caption{Влияние коэффициента обучения на траекторию ошибки traingd ($S = 10$).}
\label{fig:lr}
\end{figure}

\begin{figure}[H]
\centering
\includegraphics[width=\textwidth]{12_under_overfitting.png}
\caption{Недообучение и переобучение на зашумлённых данных.}
\label{fig:overfit}
\end{figure}

\begin{figure}[H]
\centering
\includegraphics[width=\textwidth]{13_bias_variance.png}
\caption{Ошибка на обучающих точках и относительно истинной функции в зависимости от $S$.}
\label{fig:bv}
\end{figure}

\begin{figure}[H]
\centering
\includegraphics[width=\textwidth]{00_app_screenshot.png}
\caption{Интерактивное приложение \texttt{app/nn\_trainer.html}.}
\label{fig:app}
\end{figure}
